\section{Introduction}

Microscopic leukocyte morphology remains a central source of information in hematology. Expert review of peripheral blood and bone marrow smears supports differential counts, quality control of automated analyzers, and recognition of morphologies that may warrant further laboratory interpretation. In parallel, convolutional neural networks and related image-analysis methods have achieved high accuracy on curated blood-cell image datasets, including peripheral blood and bone marrow morphology collections \cite{walter2023_ai_hematology,khan2021_wbc_review,deshpande2021_blood_cells_review,acevedo2019_peripheral_blood_cnn,matek2019_blast,matek2021_bone_marrow}. These advances make automated morphology recognition a strong research problem for biomedical image analysis. They also create a methodological risk: closed-set performance can look convincing even when the evaluation never asks what happens to cells outside the training taxonomy.

The conventional closed-set setting assumes that every test image belongs to one of the classes seen during training. That assumption is convenient for supervised learning, but it is fragile for morphology recognition. A model trained on mature leukocyte categories can encounter immature, atypical, rare, damaged, or otherwise out-of-taxonomy cells at test time. If evaluation forces every such image into one of the known classes, the model may appear confident while performing the wrong kind of task. Open-set recognition addresses this gap by requiring a system both to classify known samples and to reject samples from unseen classes \cite{scheirer2013_open_set,geng2021_osr_survey,bendale2016_openmax,yang2024_generalized_ood}. In the present study, the open-set question is deliberately narrow: can a classifier trained on a fixed leukocyte morphology taxonomy identify that a test image should not be assigned to that taxonomy?

This distinction is especially important because ``unknown'' has a technical meaning here. The known classes are basophil, eosinophil, lymphocyte, monocyte, and segmented neutrophil. An unknown sample is a morphology outside that training taxonomy. It is not a cancer label, leukemia label, malignant label, abnormal diagnosis, patient-level endpoint, or clinically positive finding. The work therefore should not be read as a diagnostic system or leukemia detector. It is an empirical study of recognition reliability when the model's label inventory is incomplete.

The hardest unknown samples are often not visually alien. A band neutrophil is close to segmented neutrophils; a monoblast can be pulled toward monocytes; blast or atypical lymphoid morphologies can resemble the neighborhood of typical lymphocytes. Such examples are not merely generic outliers. They are near-taxonomy morphologies whose errors may be structured by cell lineage, maturation, image acquisition, and representation geometry. This makes leukocyte open-set recognition a demanding biomedical image-analysis problem: the model must reject samples that are related to, but not members of, the known classes.

External domain shift compounds the problem. Hematology image datasets can differ in staining, scanner, objective, cropping, annotation policy, patient mix, image size, and class composition. Prior work in medical image analysis and blood-cell classification has shown that models trained and tested under similar imaging conditions can lose performance when evaluated across sources \cite{guan2022_domain_adaptation_medical,litjens2017_medical_ai_survey,salehi2022_cross_domain,sadafi2024_continual_wbc,tsutsui2023_wbc_domain_shift}. For open-set recognition, external validation is even more demanding because both known-class classification and unknown rejection can change simultaneously. A score that is stable internally may no longer rank methods correctly in a new acquisition domain.

Despite extensive work on leukocyte classification and growing work on open-set and out-of-distribution detection, evidence remains limited for protocols that explicitly hold out named hematological morphologies as unknown classes, evaluate multiple seeds, preserve unknowns as test-only data, and test whether internal rankings transfer to an external hematology dataset. This gap motivates the present manuscript. We evaluate a frozen open-set protocol using MLL23 as the internal domain and AML-Cytomorphology\_LMU as an independent test-only domain \cite{mll23_nature_2025,mll23_zenodo_2024,aml_lmu_tcia_2019}. The scientific thesis is that high closed-set leukocyte morphology classification performance does not imply reliable open-set recognition of unseen morphologies, particularly under morphological similarity and external domain shift.

The contributions are:
\begin{itemize}
\item A reproducible open-set evaluation protocol for leukocyte morphology recognition, with a frozen five-class known taxonomy and held-out MLL23 morphologies used only for open-set testing.
\item A five-seed and two-split sensitivity analysis showing that internal closed-set macro-F1 near 0.97 coexists with imperfect unknown rejection and high FPR@95TPR.
\item An external test-only AML-Cytomorphology\_LMU validation showing substantial domain-shift degradation and instability of internal method rankings.
\item A morphology-level failure analysis showing that difficult unknown rejection is structured by known-class attractors such as band-to-segmented neutrophil, monoblast-to-monocyte, and lymphoid-to-lymphocyte.
\item Controlled negative or secondary evidence showing that improved known-class geometry and persistent-homology descriptors do not automatically translate into robust open-set recognition.
\end{itemize}

\begin{figure}[htbp]
\centering
\includegraphics[width=\textwidth]{figure1_protocol.png}
\caption{Frozen open-set evaluation protocol. MLL23 is used for internal multiseed evaluation, with five mature leukocyte classes treated as known and 13 morphologies held out as unknown. AML-Cytomorphology\_LMU is used only as an external test domain, without training, tuning, calibration, or representation fitting.}
\label{fig:protocol}
\end{figure}
